\section{Introduction}\label{sec:intro}
% NGƯỜI VIẾT: RW (Lộc)
%
% Khung gợi ý (4 đoạn):
% 1. Bối cảnh: unit test tốn 30-50% effort; LLM hứa hẹn tự sinh test;
%    các nghiên cứu hiện có báo coverage cao nhưng chủ yếu trên hàm đơn giản.
% 2. Khoảng trống (GAP — lấy từ team-synthesis/gap-final.md):
%    - GAP-T: chưa ai đo "điểm gãy" của LLM theo cyclomatic complexity 5-10
%    - GAP-M: thiếu benchmark thống nhất branch coverage + mutation score
%      có kiểm định thống kê dưới CC được kiểm soát
% 3. Đóng góp của nghiên cứu này:
%    - Đánh giá gpt-4o-mini (one-shot) trên 120 hàm thật (60 Java + 60 Python,
%      CC 5-10) mined từ 10 repo OSS pinned commit
%    - So với 2 baseline sinh test tự động (EvoSuite, Randoop) cùng budget 60s
%    - 3 RQ với kiểm định Wilcoxon/Spearman
% 4. Cấu trúc bài báo.
%
% TODO(Lộc): viết nội dung, ~0.75-1 trang.

% RQ box — giữ nguyên wording đã được GV duyệt trong proposal:
\textbf{RQ1}: Does GPT-4o-mini achieve $\geq 80\%$ branch coverage on functions
of medium cyclomatic complexity? \\
\textbf{RQ2}: Does GPT-4o-mini achieve $\geq 60\%$ mutation score, and does it
outperform automated baselines? \\
\textbf{RQ3}: Is there a negative correlation between cyclomatic complexity and
test quality?
